% Source PDF: source/普通高考/2012/2012新课标理(河南,河北,黑龙江,吉林,海南,宁夏,山西,内蒙古,新疆,云南).pdf, page 1 (question 6).
% The source PDF contains vector line art (pdfimages -list returned no image objects).
% Node -> directed-edge list: 开始 -> 输入(N,a_1,...,a_N) ->
% k=1,A=a_1,B=a_1 -> x=a_k -> x>A; 是 -> A=x -> merge before x<B;
% 否 -> x<B; 是 -> B=x -> the same merge; 否 -> the same merge -> k>=N;
% 是 -> 输出 A,B -> 结束; 否 -> k=k+1 -> loop-back to the stem above x=a_k.
% solid segments: all node outlines and connectors; dashed segments: none.
% labels: each 是/否 sits beside its corresponding decision exit; all node text
% is locally zero-indented and horizontally/vertically centered.

\begin{tikzpicture}[
  x=1cm,
  y=1cm,
  >=Stealth,
  line width=0.45pt,
  line cap=round,
  line join=round,
  font=\normalsize,
  flowtext/.style={anchor=center, align=center,
    execute at begin node={\setlength{\parindent}{0pt}}},
  flowbox/.style={draw, minimum height=0.72cm, inner xsep=4pt, inner ysep=2pt,
    align=center, execute at begin node={\setlength{\parindent}{0pt}}},
  terminal/.style={flowbox, rounded corners=0.18cm, minimum width=1.30cm},
  io/.style={flowbox, trapezium, trapezium left angle=72, trapezium right angle=108,
    minimum height=0.78cm},
  decision/.style={flowbox, diamond, aspect=2.65, minimum width=4.20cm,
    minimum height=1.00cm, inner xsep=1pt, inner ysep=1pt},
  branchlabel/.style={flowtext, inner sep=0pt},
]
  % Main sequence and the two nested comparisons.
  \node[terminal] (start) at (0,15.5) {开始};
  \node[io, minimum width=6.30cm] (input) at (0,14.15)
    {输入 $N,a_1,a_2,\cdots,a_N$};
  \node[flowbox, minimum width=6.20cm] (init) at (0,12.75)
    {$k=1,A=a_1,B=a_1$};
  \node[flowbox, minimum width=2.55cm] (sample) at (0,11.10) {$x=a_k$};
  \node[decision] (maxcheck) at (0,9.35) {$x>A$};
  \node[flowbox, minimum width=2.55cm] (setmax) at (3.50,8.00) {$A=x$};
  \node[decision] (mincheck) at (0,6.85) {$x<B$};
  \node[flowbox, minimum width=2.55cm] (setmin) at (-4.00,5.65) {$B=x$};
  \node[decision] (countcheck) at (0,3.40) {$k\geq N$};
  \node[io, minimum width=3.10cm] (output) at (0,1.35) {输出 $A,B$};
  \node[terminal] (finish) at (0,-0.25) {结束};
  \node[flowbox, minimum width=3.00cm] (increment) at (5.35,5.70) {$k=k+1$};

  \draw[->] (start.south) -- (input.north);
  \draw[->] (input.south) -- (init.north);

  % Initialization and loop-back share a stem above x=a_k.
  \coordinate (loop-merge) at (0,11.95);
  \draw (init.south) -- (loop-merge);
  \draw[->] (loop-merge) -- (sample.north);
  \draw[->] (sample.south) -- (maxcheck.north);

  % First comparison: update A on the yes branch, otherwise test B.
  \draw[->] (maxcheck.east) -- (3.50,9.35) -- (3.50,8.48) -- (setmax.north);
  \draw[->] (maxcheck.south) -- (mincheck.north);

  % Second comparison and the explicit merge before k>=N.
  \draw[->] (mincheck.west) -- (-4.00,6.85) -- (setmin.north);
  \coordinate (merge) at (0,4.55);
  \draw (mincheck.south) -- (0,5.15);
  \draw (0,5.15) -- (merge);
  \draw (setmax.south) -- (3.50,4.55);
  \draw (3.50,4.55) -- (merge);
  \draw (setmin.south) -- (-4.00,4.55);
  \draw (-4.00,4.55) -- (merge);
  \draw[->] (merge) -- (countcheck.north);

  % Final test: output on yes, increment and loop-back on no.
  \draw[->] (countcheck.south) -- (output.north);
  \draw[->] (output.south) -- (finish.north);
  \draw[->] (countcheck.east) -- (5.35,3.40) -- (5.35,5.18) -- (increment.south);
  \draw[->] (increment.north) -- (5.35,11.95) -- (0.40,11.95);
  \draw (0.40,11.95) -- (loop-merge);

  \node[branchlabel, anchor=west] at (1.85,9.62) {是};
  \node[branchlabel, anchor=west] at (0.45,8.30) {否};
  \node[branchlabel, anchor=east] at (-2.10,7.28) {是};
  \node[branchlabel, anchor=west] at (0.45,5.95) {否};
  \node[branchlabel, anchor=west] at (1.80,3.82) {否};
\end{tikzpicture}
